\section{Introduction}

Aerial tiny-object detection places discrete training decisions under unusually coarse geometric resolution. Dense candidates lie on stride-dependent lattices, while many NMS-free detectors use a one-to-one (O2O) branch that retains a single post-conflict assigned candidate for each ground-truth object~\cite{wang2021defcn,sun2021onenet,wang2024yolov10,jocher2026yolo26}. Even a one-pixel centre displacement can change the frozen assignment winner. A small displacement can therefore alter not only overlap or score, but also the identity, rank, eligibility, or positive-set membership used to construct supervision. These are properties of the native assignment state rather than postprocessed detections.

Existing tiny-object methods modify localisation similarity, receptive-field-aware assignment, scale-aware labels, feature resolution, or supervision density~\cite{wang2021nwd,xu2022rfla,shi2024simd,xu2023dcfl,lin2026ngla,zhang2026vala,tian2026nbddla}. Dense-detector assignment methods likewise change candidate selection through adaptive matching, optimal transport, task alignment, or differentiable assignment~\cite{zhang2020atss,kim2020paa,ge2021ota,feng2021tood,zhang2019freeanchor,zhu2020autoassign}. Such methods address how positives should be constructed. They do not, by themselves, measure how stable the resolved native assignment state is after training, nor whether an apparent scale response depends on the geometric stress used to probe it.

The fact that one input pixel represents a larger fractional displacement for a smaller object is elementary. The unresolved questions are whether an apparent assignment-state scale effect survives an equivalent-area-side stress, which native states change, and whether a Top-1 exchange also implies turnover of the positive supervision set.

Measuring this state is non-trivial for three reasons. First, the instrument must reproduce native eligibility, Top-$k$ selection, conflict resolution, tie handling, and numerical precision closely enough that a reported identity change is not a tracing artefact. Second, O2O identity, O2M rank, and O2M assigned-positive-set membership are different states; for the diagnostic questions considered here, a single endpoint candidate-change variable is insufficient to distinguish the audited native states. Third, base-plus-four-direction replay observes responses only at one displacement magnitude, but does not show how close the first boundary lies, whether a trajectory is right-censored, or which native stage produced the first divergence.

We therefore develop a qualified read-only measurement of the frozen native assigner. The focal ground-truth centre is replayed under a fixed-pixel or equivalent-area-side displacement while image pixels, extracted features, model parameters, candidate grids, box size, and all non-focal ground truths remain fixed. The measurement contract is explicit: a one-pixel shift asks about sensitivity to an absolute input displacement, whereas equivalent-area-side replay scales displacement by $\sqrt{wh}$ and is not the same axis-specific fraction of width and height for a rectangular object. Neither contract is treated as a correction for the other. Analytical oracles, property invariants, mechanism mutations, stage-wise native/reference parity, exact full-grid focal-row/native equivalence, and an explicit grid-resolution gate qualify the implementation before scientific interpretation. The audit then combines endpoint replay, eligibility locking, first-divergence and multi-stage anatomy, O2O pre-resolution controls, O2M rank--set decomposition, and cardinal-direction boundary analysis. Weighted Kaplan--Meier and Aalen--Johansen summaries retain censoring and competing boundary causes.

The contributions are threefold:
\begin{itemize}[leftmargin=*,nosep]
    \item \textbf{Stress-contract dependence.}\\[-2pt] On AI-TOD-v2, fixed-pixel and equivalent-area-side replay led to qualitatively different scale conclusions.\\[-2pt] Sensitivity analyses across four protocol-eligible seeds, VisDrone, archived stages, and a compatible detector contract showed that the branch-wise response structure was not confined to a single checkpoint or analysis setting.
    \item \textbf{State-specific assignment anatomy.} Eligibility locking, first-divergence analysis, and O2M Rank--Set decomposition separate eligibility transitions, within-set rank exchange, Top-1 turnover, and positive-set turnover. For the diagnostic questions considered here, a single endpoint candidate-change variable is insufficient to distinguish the audited native states.
    \item \textbf{Cardinal boundary persistence.} Exhaustive $1/1024$ replay along four cardinal rays quantifies first-departure distance, censoring, competing boundary causes, and restricted mean persistence, revealing different observed persistence structures for assigned O2O identity and O2M assigned-set Top-1.
\end{itemize}

Implementation qualification is reported separately as a prerequisite for interpreting these measurements.

The practical role of the audit is not to prescribe a replacement assigner, but to determine whether an apparent assignment effect reflects the stress coordinate, a change in candidate ranking, a change in the supervision set, or a shift in the observed assignment boundary. Accordingly, a fixed-pixel scale effect can reflect the chosen stress coordinate rather than an intrinsic scale law; O2M Top-1 turnover need not imply positive-set turnover; and endpoint flip/no-flip omits distance, censoring, cause, and return information.
